\section{Shared-policy comparison: protocol}
\label{app:shared-cnn}

This is the protocol behind Section~\ref{sec:planning} and the behavioral rows of
Section~\ref{sec:representation}. It was fixed before any of its policy outcomes existed
(\texttt{research/shared\_decoded\_cnn\_protocol.md} in the repository).

\paragraph{Matrix.}
Four worlds: per-frame gray noise without wind, per-frame gray noise with wind, persistent gray
noise with wind, and the gate. For each world, the grid JEPA and the flow model each train three
policy seeds (0, 1, 2) inside imagination: 24 policies. The gate additionally receives three
decoded-real control policies per family: 6 policies. All 30 outcomes are reported. Gray-world JEPA
checkpoints are model seed 0; the gate JEPA is model seed 1, preselected on prior development
evidence that its decoded-real CNN could discover the strategy. Every flow checkpoint is the seed-0
fit. Policy seeds are not independent world-model fits.

\paragraph{Policy and optimizer.}
The policy is the unmodified pixel actor-critic of the original study: a three-block Impala CNN
on native $63\times63$ RGB, a GRU that also receives the previous action, a categorical actor and
a 64-bin histogram critic; 1,702,980 parameters. Both families use the same recurrent PPO: AdamW at
$2.5\times10^{-4}$ with weight decay $0.01$, four epochs, minibatch 128, 16-step truncated
backpropagation, $\gamma=0.99$, GAE $\lambda=0.9$, value coefficient $0.5$, entropy coefficient
$0.01$, clipping $0.2$, global gradient clipping $1.0$, batch 512 and horizon 48. The three non-gate
worlds use 50 updates (1,228,800 positions, 2,400 optimizer steps); every gate policy, imagined or
decoded-real, uses 150 updates (3,686,400 positions, 7,200 optimizer steps). The training script
refuses any other budget unless run as an explicitly non-scientific two-update smoke.

\paragraph{Observations.}
Each policy sees only its own model's decoded RGB: decoded posterior encodings of real frames
during burn-in and evaluation, decoded model predictions in imagination. No raw pixel reaches the
actor after the encoder. Training starts are sampled uniformly from valid eight-frame archive
histories with no reset in the history; burn-in processes the first seven decoded frames with
their real incoming actions. Latents keep the native $4\times4\times32$ interface without
resizing, whitening or amplification; decoder quantization is retained.

\paragraph{Task signals.}
Flow uses its native sampled dynamics at four function evaluations and its jointly fitted reward
and termination heads. The JEPA uses its deterministic predictor, the existing 80,000-frame
decoder, and reward and termination heads of the original architecture refitted on all 500,000
training transition labels with the original 6,000-update recipe. The refit never updates encoder,
predictor or decoder. Equal label counts do not equalize supervision stage: the flow heads send
full-strength gradients into the dynamics trunk for the first half of its training and
$0.1$-scaled gradients afterwards, while JEPA world-model fitting is reward-free.

\paragraph{Evaluation.}
256 fresh reset episodes are admitted by surviving seven uniformly random actions; the admission
never looks at the policy. The same accepted prefix identities and all environment and action keys
are reused across families and policy seeds, so every contrast is paired. The actor then controls
121 steps under the 128-step cap. Survival is return above $-0.5$; walls passed counts wall exits
before or on the terminal transition. Every policy is its final checkpoint. Native-RGB trajectories
of the first three accepted episodes are stored for each policy regardless of outcome.

\paragraph{Validation before the runs.}
Independent CPU checks of the training and evaluation code; GPU codec parity between the
imagination path and the evaluation path for every checkpoint (context versus single-step
encoding, 256- and 512-frame batches, head-cache encoding); unchanged frozen-state hashes; exact
gate-state reconstruction; shared initial-context and evaluation-packet fingerprints. Two-update
smokes per family and path supplied wall-clock estimates used only for scheduling.

\paragraph{State-blind reference rows.}
Two rules that never read an observation are run through the same paired evaluator on the same
packets: uniform random actions, and strict left--right alternation, which is the trivial survival
strategy of the still worlds (Table~\ref{tab:shared-cnn-reference}). They calibrate the survival
numbers of Table~\ref{tab:shared-cnn}. Alternation survives the per-frame still world in every
episode; both rules survive no episode of either wind world or of the gate. A trained policy whose
control length falls below these values is steering into the boundary faster than chance.

\begin{table}[h]
\centering
\caption{State-blind rules on the shared evaluation packets: survival, control length and, for the
gate, walls passed, over the same 256 paired episodes as Table~\ref{tab:shared-cnn}.}
\label{tab:shared-cnn-reference}
\inputtable{shared_cnn_reference.tex}
\end{table}

\paragraph{Imagined termination probe.}
To separate what the policy learns from what its training environment shows it, each frozen model
is rolled forward for 48 steps from 256 uniformly sampled real eight-frame contexts under the
recorded actions (the data policy is uniform random), and its own termination head is read on two
inputs: the model's predicted latent, exactly as in imagination, and the encoder latent of the real
frame at the same step. The recorded termination labels give the ground truth. Nothing is trained
and the adapters are byte-identical to those of the campaign. Figure~\ref{fig:imagined-termination}
plots the implied probability that an imagined episode is still alive after $k$ steps against the
recorded survival. Per-step outputs, including latent distances between predicted and encoder
latents, are stored with the campaign results.
